\section{Mixture-of-Models as a Model Architecture}
\label{sec:formalism}

\subsection{One model beyond one checkpoint}

In conventional deployments, a model name resolves to one checkpoint or managed
endpoint. A Mixture-of-Models (MoM) instead exposes one model interface while
computation crosses model, provider, process, and device boundaries. Its
components may be open or closed, local or remote, and independently versioned
and invoked; they are typically, but need not be, independently trained.
Its model identity is defined not by shared weights but by a versioned
executable specification that fixes logical references, admissible
substitutions, composition, constraints, and the output contract.

Each published model fixes a behavior variant $b$, its default objective, and
an allowed domain $\mathcal{U}_b$ of request preferences. For an input $x$,
conversation or session history $h$, request preference $u\in\mathcal{U}_b$,
and observable environment state $e$, we define its logical specification as
\begin{equation}
  \mathcal{M}_b =
  (\mathcal{V}, \mathcal{Q}, \mathcal{U}_b, S_{\phi}, P_{\omega}, \pi_{\psi},
   \mathcal{A}, \mathcal{G}, \mathcal{K}, \Gamma).
  \label{eq:mom-object}
\end{equation}
Here $\mathcal{V}=\{v_1,\ldots,v_n\}$ contains \emph{logical model references},
not URLs; each declares an exact identity or capability requirement together
with interfaces and deployment requirements. Optional $\mathcal{Q}$ contains
declared typed auxiliary operators for retrieval, tool use, and other bounded
effects; operators have the same explicit interface, policy, and provenance
requirements as model calls, while persistent application state remains
outside the model. $\mathcal{U}_b$ contains the typed request preferences and
policy refinements admitted by behavior variant $b$; it cannot override a hard
guard or transform one published variant into another.
$S_{\phi}$ extracts typed signals such as domain, complexity, modality, risk,
and confidence; $P_{\omega}$ projects them to calibrated scores and predicates;
and $\pi_{\psi}$ chooses the next admissible action. $\mathcal{A}$ defines
action and transition semantics, $\mathcal{G}$ contains typed preconditions and
postconditions, $\mathcal{K}$ defines intermediate and final contracts, and
$\Gamma$ imposes finite bounds on events, path length, fan-out, calls, rounds,
tokens, and time. These components may be learned or authored; content-addressed
prompts and policy assets parameterize the specification. Catalogs, topologies,
prompts, thresholds, contracts, and bounds can be optimized even without
gradients through a constituent model.

A deployment binding later resolves a logical requirement to a local checkpoint
or provider model. Thus the specification can be inspected and packaged without
credentials, then rebound without editing its decision semantics
(Section~\ref{sec:artifact}). We write
$\mathcal{R}=(\mathcal{M}_b,\mathcal{B},L)$ for a \emph{deployed MoM}, where
$\mathcal{B}$ is the environment binding and $L$ is the completed resolution
lock, including the engine revision. Logical identity belongs to $\mathcal{M}_b$.
The caller observes $\mathcal{R}$ as one virtual model: the engine revision
recorded in $L$ executes $\mathcal{M}_b$ against the resolved pool, producing an
attributed event DAG for each request. This distinction permits cross-engine
portability without reducing the deployed model to either the engine or its
endpoints.

\subsection{Execution-trace semantics}

Before event $j$, the runtime state is
\begin{equation}
  \xi_j=(x,h,u,e,z_j,m_j,b_j), \qquad
  z_j=P_{\omega}(S_{\phi}(x,h,e,m_j),h,u,e,m_j),
\end{equation}
where $m_j$ is typed working memory and $b_j$ is the remaining resource budget.
Recomputing $z_j$ allows later decisions to use validated outputs from earlier
events without granting a planner unbounded control.
The primitive actions are
\begin{equation}
\begin{split}
  \textsc{Call}(v,r),\quad &\textsc{Op}(q,r),\quad
  \textsc{Fork}(a_1,\ldots,a_k),\\
  \textsc{Join}(f),\quad &\textsc{Check}(k),\quad
  \textsc{Transform}(f),\quad
  \textsc{Return}(s,y).
\end{split}
\label{eq:actions}
\end{equation}
A call invokes $v$ with a checked request; $\textsc{Op}$ invokes a declared,
typed external operator $q$ such as retrieval or a tool; fork and join implement
bounded parallel composition; check evaluates a guard; and transform is typed
and deterministic.  Planners and judges are ordinary calls, not sources of
unbounded control flow. Every terminal trace ends with a typed status
$s\in\{\texttt{ok},\texttt{abstain},\texttt{reject},\texttt{error}\}$; for a
non-answer status, $y=\bot$. Let $\bar y=(s,y)$ denote the complete terminal
outcome.

An execution trace is a finite, labeled event DAG
\begin{equation}
  \tau=(N,\prec,\sigma,\{(\xi_j,a_j,o_j)\}_{j\in N},\bar y),
  \label{eq:trace}
\end{equation}
where $N$ is the event set, $\prec$ is the control-and-data happens-before
relation, and $\sigma$ is a canonical topological order used only for logging
and replay. A sequential trace has a total order; \textsc{Fork} creates
incomparable child regions, and \textsc{Join} depends on their terminal events.
Stable branch identifiers make $\sigma$ unique, so different interleavings do
not create duplicate traces. The state $\xi_j$ is a deterministic function of
the outcomes visible through predecessors of $j$; wall-clock observations and
timeouts remain part of $e$ and $o_j$. Both $|N|$ and every path length are
finite under $\Gamma$.

The policy may be stochastic and condition on prior model outputs through
$m_j$. Let $L$ identify the bound constituent revisions, let $C(\tau)\subseteq
N$ contain the controller's choice events, and let $E(\tau)\subseteq N$ contain
effectful model and operator calls. Graph-prescribed and administrative events
have unit probability. For event $j$, the causal history
$H_j=\{o_k:k\prec j\}$ contains only predecessor-visible outcomes; $\xi_j$ is
measurable with respect to this history. In particular, a branch cannot observe
an incomparable branch before their declared join.

Parallel calls can share provider, load, and runtime randomness, so we do not
assume that their outcomes are conditionally independent. Instead, let
$K_L$ be the causal execution kernel over all effectful outcomes for the chosen
action DAG, and let $r_{E(\tau)}$ collect the checked requests attached to its
effectful actions. The joint probability of a valid trace and terminal outcome is
\begin{equation}
\begin{aligned}
 p_{\mathcal{R}}(\bar y,\tau\mid x,h,u,e)
 ={}& \mathbf{1}[\textsc{Return}(\tau)=\bar y]
      \prod_{j\in C(\tau)}\pi_{\psi}(a_j\mid\xi_j)\\
 &\cdot K_L\!\left(o_{E(\tau)}\mid r_{E(\tau)},a_{E(\tau)},
                   \prec,x,h,u,e\right).
\end{aligned}
\label{eq:trace-factorization}
\end{equation}
The kernel is causal: an event distribution may depend on $H_j$ and shared
runtime state, but not on incomparable or future outcomes. When effectful calls
are conditionally independent given their causal histories, $K_L$ reduces to a
product of per-call response kernels. The user-facing distribution of the
deployed realization marginalizes admissible event DAGs,
\begin{equation}
 p_{\mathcal{R}}(\bar y\mid x,h,u,e)
 = \sum_{\tau\in\mathcal{T}_{\mathrm{adm}}}
   p_{\mathcal{R}}(\bar y,\tau\mid x,h,u,e).
 \label{eq:mom-distribution}
\end{equation}
Closed-model token probabilities need not be accessible: training may use
sampled outputs, scores, preferences, or outcomes.

The trace view unifies common structures.  \textbf{Selection} calls one model;
a \textbf{cascade} conditionally escalates through an ordered sequence;
\textbf{fusion} forks a bounded panel and joins its responses;
\textbf{multi-round collaboration} repeats calls over typed shared state; and a
\textbf{workflow} instantiates a bounded static graph or validated planner
output.  Routing is therefore one MoM topology, while a collection of APIs is
not a MoM until its admissible traces, contracts, and identity are explicit.
Fixed panels and workflows are degenerate policies that assign probability one
to one graph; conditional computation does not require every trace to activate
a strict subset of the catalog.

\begin{figure}[H]
  \centering
  \resizebox{\linewidth}{!}{\begin{tikzpicture}[
  x=1cm,y=1cm,font=\sffamily,>=Latex,
  edge/.style={-{Latex[length=1.6mm]},draw=MoMGray,line width=0.55pt},
  call/.style={draw=MoMBorder,fill=white,rounded corners=1.5pt,line width=0.5pt,minimum width=0.62cm,minimum height=0.45cm,font=\tiny},
  selected/.style={draw=MoMDark,fill=MoMBlue,rounded corners=1.5pt,line width=0.7pt,minimum width=0.62cm,minimum height=0.45cm,font=\tiny},
  control/.style={draw=none,fill=MoMDark,text=white,rounded corners=1.5pt,minimum width=0.62cm,minimum height=0.45cm,font=\tiny},
  panel/.style={draw=MoMBorder,rounded corners=3pt,line width=0.55pt},
  title/.style={font=\sffamily\bfseries\fontsize{7.0}{8}\selectfont,text=MoMText},
  note/.style={font=\sffamily\fontsize{5.8}{6.7}\selectfont,text=MoMMuted}
]
\foreach \x/\name in {0/Selection,3.28/Cascade,6.56/Fusion,9.84/Multi-round,13.12/Workflow}{
  \draw[panel] (\x,0) rectangle ++(3.04,2.55);
  \node[title] at (\x+1.52,2.27) {\name};
}

% Selection.
\node[control] (s0) at (0.52,1.22) {$\pi$};
\node[selected] (s1) at (1.56,1.62) {$v_1$};
\node[call] (s2) at (1.56,0.82) {$v_2$};
\node[call] (s3) at (2.52,1.22) {$y$};
\draw[edge] (s0) -- (s1);
\draw[draw=MoMBorder,line width=0.45pt] (s0) -- (s2);
\draw[edge] (s1) -- (s3);
\node[note] at (1.52,0.28) {one active call};

% Cascade.
\node[selected] (c1) at (3.72,1.22) {$v_1$};
\node[control] (ck) at (4.58,1.22) {?};
\node[selected] (c2) at (5.45,1.22) {$v_2$};
\node[call] (cy) at (6.17,1.22) {$y$};
\draw[edge] (c1) -- (ck);
\draw[edge] (ck) -- (c2);
\draw[edge] (c2) -- (cy);
\node[note] at (4.80,0.28) {conditional escalation};

% Fusion.
\node[control] (f0) at (6.98,1.22) {fork};
\node[selected] (f1) at (8.08,1.65) {$v_1$};
\node[selected] (f2) at (8.08,0.80) {$v_2$};
\node[control] (fj) at (9.16,1.22) {join};
\draw[edge] (f0) -- (f1);
\draw[edge] (f0) -- (f2);
\draw[edge] (f1) -- (fj);
\draw[edge] (f2) -- (fj);
\node[note] at (8.08,0.28) {bounded parallel panel};

% Multi-round.
\node[selected] (r1) at (10.24,1.22) {$v_1$};
\node[selected] (r2) at (11.30,1.22) {$v_2$};
\node[control] (rb) at (12.36,1.22) {$\leq R$};
\draw[edge] (r1) -- (r2);
\draw[edge] (r2) -- (rb);
\draw[edge] (rb.south) to[bend left=35] (r1.south);
\node[note] at (11.36,0.28) {typed shared state};

% Workflow.
\node[control] (w0) at (13.50,1.55) {plan};
\node[selected] (w1) at (14.54,1.82) {$v_1$};
\node[selected] (w2) at (14.54,1.04) {$v_2$};
\node[call] (w3) at (15.63,1.43) {tool};
\draw[edge] (w0) -- (w1);
\draw[edge] (w0) -- (w2);
\draw[edge] (w1) -- (w3);
\draw[edge] (w2) -- (w3);
\node[note] at (14.64,0.28) {validated finite graph};
\end{tikzpicture}
}
  \caption{\textbf{One trace semantics, multiple execution topologies.}
  Selection, conditional cascades, bounded parallel fusion, multi-round
  collaboration, and validated workflows use the same primitive actions and
  resource bounds.  Blue nodes are calls realized in the illustrated trace;
  white nodes remain admissible but inactive.}
  \label{fig:trace-topologies}
\end{figure}

\paragraph{Coverage of bounded graphs.}
The action vocabulary is designed to represent finite, well-typed control-flow graphs
whose effectful vertices call members of $\mathcal{V}\cup\mathcal{Q}$, whose
branches and joins have declared semantics, and whose cycles have finite bounds.
Model and operator vertices map to \textsc{Call} and \textsc{Op}; branches to
\textsc{Check} and policy transitions; parallel regions to
\textsc{Fork}/\textsc{Join}; deterministic vertices to \textsc{Transform}; and
terminals to \textsc{Return}. Bounded cycles can be unrolled or represented by
a counter in typed memory. For fixed binding, response kernels, and scheduler
semantics, the translation preserves typed terminal outcomes and effectful
calls up to deterministic administrative steps under these declared
assumptions. This is a design property, not a claim that every graph can be
learned efficiently. It lets
selection, cascades, panels, multi-round protocols, and finite workflows share
one runtime action vocabulary.

\subsection{Preferences define objectives; policies must satisfy constraints}
\label{sec:objectives}

Let $R(\bar y,y^*)$ measure task utility, including the utility of explicit
abstention, rejection, or error, and let
$c(\tau,e)\in\mathbb{R}_{+}^{d}$ contain soft resource outcomes such as monetary
cost, latency, generated tokens, energy use, and estimated emissions. A
request preference $u$ compiles to weights $\lambda(u)$ and, where appropriate,
targets or service-level objectives.  For a data distribution $\mathcal{D}$, a
basic training objective is
\begin{equation}
 \max_{\Theta}
 \;\mathbb{E}_{(x,h,u,e,y^*)\sim\mathcal{D},\,
 \tau\sim p_{\mathcal{R}_{\Theta}}(\cdot\mid x,h,u,e)}
 \left[R(\bar y,y^*)-\lambda(u)^{\top}c(\tau,e)\right],
 \label{eq:soft-objective}
\end{equation}
where $\Theta=(\theta_{\mathrm{log}},\beta)$ denotes the logical and deployment
coordinates introduced in Section~\ref{sec:optimization}.

Privacy, residency, model allowlists, and capability rules that are known
before a call are typed action preconditions
$g_\ell^{\mathrm{pre}}(\xi_j,a)\in\{0,1,\star\}$, where $\star$ denotes
unknown evidence. They define
\begin{equation}
 \mathcal{A}_{\mathrm{adm}}(\xi_j)
 =\{a\in\mathcal{A}(\xi_j):
 g_\ell^{\mathrm{pre}}(\xi_j,a)=1\;\forall \ell\},
 \qquad
 \pi_\psi(a\mid\xi_j)=0\ \text{if }a\notin\mathcal{A}_{\mathrm{adm}}(\xi_j).
 \label{eq:admissible}
\end{equation}
Only a verified value of $1$ admits an action. The guards cover graph
reachability, artifact policy, environment and binding policy, the selected
behavior variant, and request preference. Either may remove actions from the
environment-admissible set but cannot restore an action excluded by residency,
authorization, or deployment policy.
Postconditions capture properties---such as output validity or grounding---that
can be evaluated only after an action returns. A trace may return \texttt{ok}
only if all postconditions on its realized path hold; otherwise it must take a
declared repair, fallback, abstention, or rejection edge. Thus
$\mathcal{T}_{\mathrm{adm}}\subseteq\mathcal{T}_{\Gamma}$ contains traces whose
actions satisfy Equation~\ref{eq:admissible} and whose failed postconditions
follow declared transitions. If no action is admissible, execution fails closed
with $\textsc{Return}(\texttt{reject},\bot)$. Statistical requirements such as
tail latency should instead use population constraints such as
$\Pr[g_k^{\mathrm{post}}=0]\leq\delta_k$; an unobserved outcome cannot be
filtered before execution.
These guarantees are relative to the declared evidence and checker semantics;
they do not establish confidentiality, safety, or factuality unless those
mechanisms are themselves sound.

Published behavior variants such as \emph{cost-first}, \emph{accuracy-first},
\emph{balanced}, \emph{security-first}, and \emph{grounding-first} fix both
objective defaults and non-negotiable guards, preventing quality from
compensating for a privacy or contract violation. Request preferences may tune
declared tradeoffs only within the selected variant's domain.

Two operating points make the distinction precise. With quality floor $q$,
nonnegative resource weights $w_c$, and expected resource cap $\bar b$, the
same MoM may be optimized as
\begin{equation}
\begin{aligned}
 \text{cost-first:}\quad
 &\min_{\Theta}\ \mathbb{E}[w_c^{\top}c(\tau,e)]
 &&\text{s.t.}\quad \mathbb{E}[R(\bar y,y^*)]\ge q,\\
 \text{accuracy-first:}\quad
 &\max_{\Theta}\ \mathbb{E}[R(\bar y,y^*)]
 &&\text{s.t.}\quad \mathbb{E}[c(\tau,e)]\preceq \bar b.
\end{aligned}
\label{eq:dual-profiles}
\end{equation}
The first uses conditional computation to avoid unnecessary work; the second
uses an expected resource envelope to exploit complementary failure modes.
Per-trace hard limits remain enforced by $\Gamma$. These are first-class
behavior variants with separate versioned coordinates, not application-side
graphs.

\paragraph{Executor contract.}
A conforming executor must make the preceding semantics operational. Given a
well-typed specification, finite event, fan-out, path, and action-duration
bounds, every budget exhaustion, timeout, or stuck state must follow a declared
transition, and execution must eventually terminate within $\Gamma$ with a
typed status. The controller's support must remain inside
$\mathcal{A}_{\mathrm{adm}}(\xi_j)$, and every successful return must be guarded
by its declared postconditions. These obligations yield three checkable runtime
invariants: execution terminates within $\Gamma$ with a typed status; no action
with a false declared precondition is issued; and \texttt{ok} is reachable only
after required checks succeed. They are executor guarantees relative to the
declared guards and evidence, not statistical claims about the correctness of
those guards.

\subsection{Architectural invariants}

Equation~\ref{eq:mom-object} imposes five invariants: execution is
conditionally selected and resource-bounded;
composition is typed; recursion, iteration, and fan-out are bounded; logical
semantics are separate from placement; and every response is attributable to a
MoM version, binding, and trace. These invariants make a MoM optimizable like a
model and operable like a distributed system.
